% GENERATED FILE. Edit canonical lessons, then run npm run generate:books.
% canonical-source-hash: fnv1a64:9fdae644a90537d7
% canonical-lessons: ZH-C447-buxing, ZH-C447-changcheng, ZH-C447-buxingjie, ZH-C447-raolu, ZH-C447-renxingdao, ZH-R447-review-442-447

\chapter{On Foot}
\label{ch:zh-on-foot}
\hlchaptermodality{447}

\begin{chapteropening}
\textbf{By the end of this chapter:} \emph{I can talk about walking, the Great Wall, a pedestrian street, going the long way round and the pavement.}

Complete the last lesson of chapter 447: I can talk about walking, the Great Wall, a pedestrian street, going the long way round and the pavement.
\end{chapteropening}

\section[bùxíng]{bùxíng — to get there on foot}

\label{lesson:ZH-C447-buxing}



\begin{quote}
Before the new one: say the Mandarin for the wall round a city, then the Mandarin for to move to the city.
\end{quote}

\subsection*{You'll want to know: bùxíng}
\emph{bùxíng} — \textquotedblleft{}to get there on foot\textquotedblright{}.

\begin{quote}
\emph{cóng jiǔdiàn bùxíng qù gǔchéng yào shí fēnzhōng} — It's a ten-minute walk from the hotel to the old town.
\end{quote}

Its characters come later; for now it is a word for the ear.

\begin{cousinweb}
\emph{bùxíng} is built from \emph{bù} 'a step' and \emph{xíng} 'to go', both inherited from Middle Chinese, \emph{buH} and \emph{haeng} (Guangyun, in Baxter's transcription), which glosses the first as a walking step and the second 'to walk'.
\end{cousinweb}

\begin{grammarlens}[title={What you've built}]
One more word for getting around.
\end{grammarlens}

\subsection*{Your turn}
\begin{itemize}
\raggedright
  \item \emph{Say it:} \emph{bùxíng}
  \item \emph{Say it:} \emph{bùxíng}, once more
  \item \emph{Say it:} \emph{cóng jiǔdiàn bùxíng qù gǔchéng yào shí fēnzhōng} — It's a ten-minute walk from the hotel to the old town.
  \item \emph{Recall:} say the Mandarin for the wall round a city, then the Mandarin for to move to the city, then say \emph{bùxíng} again
\end{itemize}

\begin{tcolorbox}[breakable,colback=teal!4,colframe=teal!35!black,title={Before you move on}]
What does bùxíng mean? (\textquotedblleft{}To get there on foot\textquotedblright{}.) Say it once more.
\end{tcolorbox}
\section[\texorpdfstring{\zh{长城} (Chángchéng)}{长城 (Chángchéng)}]{\zh{长城} (Chángchéng) — the Great Wall}

\label{lesson:ZH-C447-changcheng}



\begin{quote}
Before the new one: say the Mandarin for to move to the city, then the Mandarin for to get there on foot.
\end{quote}

\subsection*{You'll want to know: \zh{长城}}
\textbf{\zh{长城}} — \emph{Chángchéng} — \textquotedblleft{}the Great Wall\textquotedblright{}.

\begin{quote}
\textbf{\zh{我们坐车去长城}} \emph{wǒmen zuò chē qù Chángchéng} — We're going to the great wall by bus.
\end{quote}

\begin{cousinweb}
\textbf{\zh{长城}} (\emph{Chángchéng}) is built from \zh{长} 'long' and \zh{城} 'a wall round a city', both inherited from Middle Chinese, \emph{drjang} and \emph{dzyeng} (Guangyun, in Baxter's transcription), which glosses the first 'long' and 'far'.
\end{cousinweb}

\begin{grammarlens}[title={What you've built}]
One more word for getting around.
\end{grammarlens}

\subsection*{Your turn}
\begin{itemize}
\raggedright
  \item \emph{Say it:} \emph{Chángchéng}
  \item \emph{Say it:} \emph{Chángchéng}, once more
  \item \emph{Say it:} \emph{wǒmen zuò chē qù Chángchéng} — We're going to the great wall by bus.
  \item \emph{Recall:} say the Mandarin for to move to the city, then the Mandarin for to get there on foot, then say \emph{Chángchéng} again
\end{itemize}

\begin{tcolorbox}[breakable,colback=teal!4,colframe=teal!35!black,title={Before you move on}]
What does \zh{长城} mean? (\textquotedblleft{}The Great Wall\textquotedblright{}.) Say it once more.
\end{tcolorbox}
\section[bùxíngjiē]{bùxíngjiē — pedestrian street; a street closed to cars}

\label{lesson:ZH-C447-buxingjie}



\begin{quote}
Before the new one: say the Mandarin for to get there on foot, then the Mandarin for the Great Wall.
\end{quote}

\subsection*{You'll want to know: bùxíngjiē}
\emph{bùxíngjiē} — \textquotedblleft{}pedestrian street; a street closed to cars\textquotedblright{}.

\begin{quote}
\emph{qù bùxíngjiē zěnme zǒu} — How do I get to the pedestrian street?
\end{quote}

Its characters come later; for now it is a word for the ear.

\begin{cousinweb}
\emph{bùxíngjiē} is built from \emph{bùxíng} 'on foot' and \emph{jiē} 'a street', all inherited from Middle Chinese, \emph{buH}, \emph{haeng} and \emph{kea} (Guangyun, in Baxter's transcription), which glosses the last as a road that runs four ways.
\end{cousinweb}

\begin{grammarlens}[title={What you've built}]
One more word for getting around.
\end{grammarlens}

\subsection*{Your turn}
\begin{itemize}
\raggedright
  \item \emph{Say it:} \emph{bùxíngjiē}
  \item \emph{Say it:} \emph{bùxíngjiē}, once more
  \item \emph{Say it:} \emph{qù bùxíngjiē zěnme zǒu} — How do I get to the pedestrian street?
  \item \emph{Recall:} say the Mandarin for to get there on foot, then the Mandarin for the Great Wall, then say \emph{bùxíngjiē} again
\end{itemize}

\begin{tcolorbox}[breakable,colback=teal!4,colframe=teal!35!black,title={Before you move on}]
What does bùxíngjiē mean? (\textquotedblleft{}Pedestrian street; a street closed to cars\textquotedblright{}.) Say it once more.
\end{tcolorbox}
\section[ràolù]{ràolù — to take a detour; to take the long way round}

\label{lesson:ZH-C447-raolu}



\begin{quote}
Before the new one: say the Mandarin for the Great Wall, then the Mandarin for a street closed to cars.
\end{quote}

\subsection*{You'll want to know: ràolù}
\emph{ràolù} — \textquotedblleft{}to take a detour; to take the long way round\textquotedblright{}.

\begin{quote}
\emph{qiánbian zài xiū lù wǒmen děi ràolù} — There are roadworks ahead, so we'll have to go the long way round.
\end{quote}

Its characters come later; for now it is a word for the ear.

\begin{cousinweb}
\emph{ràolù} is built from \emph{rào} 'to wind round' and \zh{路} 'a road', both inherited from Middle Chinese, \emph{nyewX} and \emph{luH} (Guangyun, in Baxter's transcription), which glosses the first 'to wind round'; today's falling tone is irregular beside it.
\end{cousinweb}

\begin{grammarlens}[title={What you've built}]
One more word for getting around.
\end{grammarlens}

\subsection*{Your turn}
\begin{itemize}
\raggedright
  \item \emph{Say it:} \emph{ràolù}
  \item \emph{Say it:} \emph{ràolù}, once more
  \item \emph{Say it:} \emph{qiánbian zài xiū lù wǒmen děi ràolù} — There are roadworks ahead, so we'll have to go the long way round.
  \item \emph{Recall:} say the Mandarin for the Great Wall, then the Mandarin for a street closed to cars, then say \emph{ràolù} again
\end{itemize}

\begin{tcolorbox}[breakable,colback=teal!4,colframe=teal!35!black,title={Before you move on}]
What does ràolù mean? (\textquotedblleft{}To take a detour; to take the long way round\textquotedblright{}.) Say it once more.
\end{tcolorbox}
\section[rénxíngdào]{rénxíngdào — pavement; sidewalk}

\label{lesson:ZH-C447-renxingdao}



\begin{quote}
Before the new one: say the Mandarin for a street closed to cars, then the Mandarin for to take the long way round.
\end{quote}

\subsection*{You'll want to know: rénxíngdào}
\emph{rénxíngdào} — \textquotedblleft{}pavement; sidewalk\textquotedblright{}.

\begin{quote}
\emph{qǐng zǒu rénxíngdào bié zǒu mǎlù} — Please keep to the pavement and stay off the road.
\end{quote}

Its characters come later; for now it is a word for the ear.

\begin{cousinweb}
\emph{rénxíngdào} is built from \zh{人} 'people', \emph{xíng} 'to walk' and \zh{道} 'a way', all inherited from Middle Chinese, \emph{nyin}, \emph{haeng} and \emph{dawX} (Guangyun, in Baxter's transcription).
\end{cousinweb}

\begin{grammarlens}[title={What you've built}]
One more word for getting around.
\end{grammarlens}

\subsection*{Your turn}
\begin{itemize}
\raggedright
  \item \emph{Say it:} \emph{rénxíngdào}
  \item \emph{Say it:} \emph{rénxíngdào}, once more
  \item \emph{Say it:} \emph{qǐng zǒu rénxíngdào bié zǒu mǎlù} — Please keep to the pavement and stay off the road.
  \item \emph{Recall:} say the Mandarin for a street closed to cars, then the Mandarin for to take the long way round, then say \emph{rénxíngdào} again
\end{itemize}

\begin{tcolorbox}[breakable,colback=teal!4,colframe=teal!35!black,title={Before you move on}]
What does rénxíngdào mean? (\textquotedblleft{}Pavement; sidewalk\textquotedblright{}.) Say it once more.
\end{tcolorbox}
\section[(dialogue)]{Looking back: from The River to On Foot}

\label{lesson:ZH-R447-review-442-447}



\begin{quote}
Say the Mandarin for the far bank and a river.
\end{quote}

\subsection*{Your turn}
\begin{itemize}
\raggedright
  \item \emph{Say it:} the far bank, a river, upstream, to rise after rain, a great river
  \item \emph{Say it:} crystal-clear, to keep from harm, a nature reserve, to escort someone, to catch fish
  \item \emph{Say it:} to dive under water, an island, a peninsula, a lighthouse, sunscreen
  \item \emph{Say it:} a wave of the sea, a fishing boat, to go ashore, a seashell, to ebb
  \item \emph{Say it:} a narrow old lane, a walled town, an ancient walled town, the wall round a city, to move to the city
  \item \emph{Say it:} to get there on foot, the Great Wall, a street closed to cars, to take the long way round, the pavement beside the road
  \item \emph{Read it:} \zh{河} \zh{护} \zh{岛} \zh{城} — each one aloud, then a word that holds it
\end{itemize}

\begin{tcolorbox}[breakable,colback=teal!4,colframe=teal!35!black,title={Before you move on}]
The far bank? (\emph{duì'àn}.) A river? (\textbf{\zh{河}}, \emph{hé}.) Upstream? (\textbf{\zh{上游}}, \emph{shàngyóu}.) To rise after rain? (\emph{zhǎngshuǐ}.) A great river? (\emph{jiāng}.) Crystal-clear? (\emph{qīngchè}.) To keep from harm? (\textbf{\zh{保护}}, \emph{bǎohù}.) A nature reserve? (\textbf{\zh{保护区}}, \emph{bǎohùqū}.) To escort someone? (\textbf{\zh{护送}}, \emph{hùsòng}.) To catch fish? (\emph{bǔyú}.) To dive under water? (\emph{qiánshuǐ}.) An island? (\textbf{\zh{岛}}, \emph{dǎo}.) A peninsula? (\textbf{\zh{半岛}}, \emph{bàndǎo}.) A lighthouse? (\emph{dēngtǎ}.) Sunscreen? (\emph{fángshàishuāng}.) A wave of the sea? (\emph{hǎilàng}.) A fishing boat? (\emph{yúchuán}.) To go ashore? (\emph{shàng'àn}.) A seashell? (\emph{bèiké}.) To ebb? (\emph{tuìcháo}.) A narrow old lane? (\emph{hútòng}.) A walled town? (\textbf{\zh{城}}, \emph{chéng}.) An ancient walled town? (\textbf{\zh{古城}}, \emph{gǔchéng}.) The wall round a city? (\emph{chéngqiáng}.) To move to the city? (\textbf{\zh{进城}}, \emph{jìnchéng}.) To get there on foot? (\emph{bùxíng}.) The Great Wall? (\textbf{\zh{长城}}, \emph{Chángchéng}.) A street closed to cars? (\emph{bùxíngjiē}.) To take the long way round? (\emph{ràolù}.) The pavement beside the road? (\emph{rénxíngdào}.)
\end{tcolorbox}
